\documentclass[12pt,a4paper]{article}
\usepackage[T1]{fontenc}
\usepackage[utf8]{inputenc}
\usepackage[expansion=false]{microtype}
\usepackage[margin=2.5cm]{geometry}
\usepackage{amsmath,amssymb}
\usepackage{booktabs,array,graphicx}
\usepackage{xcolor}
\usepackage{setspace}
\usepackage[numbers,sort&compress]{natbib}
\usepackage{algorithm}
\usepackage{algpseudocode}
\usepackage[colorlinks=true,linkcolor=blue!50!black,citecolor=blue!50!black,urlcolor=blue!50!black]{hyperref}
\setstretch{1.15}

% ---- draft switch: \drafttrue shows red TODO notes; set \draftfalse once every item is resolved --------------------
\newif\ifdraft \drafttrue
\newcommand{\todo}[1]{\ifdraft{\color{red}\textbf{[TODO: #1]}}\fi}

% ---- numbers generated by the code (one run_id for every table, macro and statement) -----------------------------
% generated by run_experiments.py -- run_id ffb593314ae3; protocol_conformant=True
\newcommand{\RunID}{ffb593314ae3}
\newcommand{\RegAdaptive}{0.0202}
\newcommand{\RegLoAdaptive}{0.0189}
\newcommand{\RegHiAdaptive}{0.0214}
\newcommand{\RegSDAdaptive}{0.0064}
\newcommand{\RegAHPTen}{0.0002}
\newcommand{\RegLoAHPTen}{0.0001}
\newcommand{\RegHiAHPTen}{0.0002}
\newcommand{\RegSDAHPTen}{0.0002}
\newcommand{\RegAHPStatic}{0.0945}
\newcommand{\RegLoAHPStatic}{0.0917}
\newcommand{\RegHiAHPStatic}{0.0974}
\newcommand{\RegSDAHPStatic}{0.0146}
\newcommand{\PreAdaptive}{0.0038}
\newcommand{\PostAdaptive}{0.0366}
\newcommand{\PreAHPTen}{0.0001}
\newcommand{\PostAHPTen}{0.0002}
\newcommand{\TransAdaptive}{0.0725}
\newcommand{\LateAdaptive}{0.0007}
\newcommand{\DiffAHPTen}{$-0.0200$}
\newcommand{\DiffAHPTenLo}{$-$0.0213}
\newcommand{\DiffAHPTenHi}{$-$0.0188}
\newcommand{\AddHolm}{0.025}
\newcommand{\AddDiff}{$-0.0005$}
\newcommand{\AddSig}{yes}
\newcommand{\LOneZero}{1.0}
\newcommand{\OracleMultRegret}{0.0004}
\newcommand{\OracleAddRegret}{0.0000}
\newcommand{\CRMatrices}{600}
\newcommand{\CRMean}{0.0036}
\newcommand{\CRSD}{0.0006}
\newcommand{\CRMedian}{0.0035}
\newcommand{\CRQOne}{0.0031}
\newcommand{\CRQThree}{0.0040}
\newcommand{\CRProp}{1.000}
\newcommand{\ClipN}{1{,}968}
\newcommand{\ClipTotal}{9{,}000}
\newcommand{\ClipPct}{21.9}
\newcommand{\NoisePSDFZero}{0.0198}
\newcommand{\NoisePSDFHigh}{0.0238}
\newcommand{\SynPosPointZero}{1}
\newcommand{\SynPosSigZero}{1}
\newcommand{\SynNegSigZero}{4}
\newcommand{\SynNegPointZero}{5}
\newcommand{\SynPosPointHalf}{2}
\newcommand{\SynPosSigHalf}{1}
\newcommand{\SynNegSigHalf}{4}
\newcommand{\SynNegPointHalf}{4}
\newcommand{\SynPosPointOne}{6}
\newcommand{\SynPosSigOne}{4}
\newcommand{\SynNegSigOne}{0}
\newcommand{\SynNegPointOne}{0}
\newcommand{\SynPosPointOneHalf}{6}
\newcommand{\SynPosSigOneHalf}{6}
\newcommand{\SynNegSigOneHalf}{0}
\newcommand{\SynNegPointOneHalf}{0}
\newcommand{\SynCellLowRevAGapA}{$-0.0200$$^{*}$}
\newcommand{\SynCellLowRevAGapB}{$-0.0141$$^{*}$}
\newcommand{\SynCellLowRevAGapC}{$+0.0071$}
\newcommand{\SynCellLowRevAGapD}{$+0.0314$$^{*}$}
\newcommand{\SynCellLowRevBGapA}{$+0.0022$$^{*}$}
\newcommand{\SynCellLowRevBGapB}{$+0.0065$$^{*}$}
\newcommand{\SynCellLowRevBGapC}{$+0.0236$$^{*}$}
\newcommand{\SynCellLowRevBGapD}{$+0.0447$$^{*}$}
\newcommand{\SynCellLowRevCGapA}{$-0.0108$$^{*}$}
\newcommand{\SynCellLowRevCGapB}{$-0.0055$$^{*}$}
\newcommand{\SynCellLowRevCGapC}{$+0.0139$$^{*}$}
\newcommand{\SynCellLowRevCGapD}{$+0.0368$$^{*}$}
\newcommand{\SynCellHighRevAGapA}{$-0.0237$$^{*}$}
\newcommand{\SynCellHighRevAGapB}{$-0.0177$$^{*}$}
\newcommand{\SynCellHighRevAGapC}{$+0.0034$}
\newcommand{\SynCellHighRevAGapD}{$+0.0278$$^{*}$}
\newcommand{\SynCellHighRevBGapA}{$-0.0016$}
\newcommand{\SynCellHighRevBGapB}{$+0.0027$}
\newcommand{\SynCellHighRevBGapC}{$+0.0198$$^{*}$}
\newcommand{\SynCellHighRevBGapD}{$+0.0409$$^{*}$}
\newcommand{\SynCellHighRevCGapA}{$-0.0145$$^{*}$}
\newcommand{\SynCellHighRevCGapB}{$-0.0092$$^{*}$}
\newcommand{\SynCellHighRevCGapC}{$+0.0102$$^{*}$}
\newcommand{\SynCellHighRevCGapD}{$+0.0332$$^{*}$}
\newcommand{\SynZeroNSCells}{t=33, SD=0.20 (-0.0016)}
\newcommand{\EWDiff}{$-0.0005$}
\newcommand{\EWLo}{$-$0.0011}
\newcommand{\EWHi}{0.0002}
\newcommand{\AHPShrinkLate}{0.0108}
\newcommand{\AHPAdvEarly}{0.0200}
\newcommand{\SemiPosPointZero}{1}
\newcommand{\SemiPosSigZero}{0}
\newcommand{\SemiNegSigZero}{5}
\newcommand{\SemiNegPointZero}{5}
\newcommand{\SemiPosPointHalf}{3}
\newcommand{\SemiPosSigHalf}{2}
\newcommand{\SemiNegSigHalf}{2}
\newcommand{\SemiNegPointHalf}{3}
\newcommand{\SemiPosPointOne}{6}
\newcommand{\SemiPosSigOne}{5}
\newcommand{\SemiNegSigOne}{0}
\newcommand{\SemiNegPointOne}{0}
\newcommand{\SemiPosPointOneHalf}{6}
\newcommand{\SemiPosSigOneHalf}{6}
\newcommand{\SemiNegSigOneHalf}{0}
\newcommand{\SemiNegPointOneHalf}{0}
\newcommand{\SemiCellLowRevAGapA}{$-0.0060$$^{*}$}
\newcommand{\SemiCellLowRevAGapB}{$-0.0024$$^{*}$}
\newcommand{\SemiCellLowRevAGapC}{$+0.0047$$^{*}$}
\newcommand{\SemiCellLowRevAGapD}{$+0.0115$$^{*}$}
\newcommand{\SemiCellLowRevBGapA}{$+0.0007$}
\newcommand{\SemiCellLowRevBGapB}{$+0.0041$$^{*}$}
\newcommand{\SemiCellLowRevBGapC}{$+0.0104$$^{*}$}
\newcommand{\SemiCellLowRevBGapD}{$+0.0158$$^{*}$}
\newcommand{\SemiCellLowRevCGapA}{$-0.0028$$^{*}$}
\newcommand{\SemiCellLowRevCGapB}{$+0.0011$}
\newcommand{\SemiCellLowRevCGapC}{$+0.0080$$^{*}$}
\newcommand{\SemiCellLowRevCGapD}{$+0.0143$$^{*}$}
\newcommand{\SemiCellHighRevAGapA}{$-0.0080$$^{*}$}
\newcommand{\SemiCellHighRevAGapB}{$-0.0044$$^{*}$}
\newcommand{\SemiCellHighRevAGapC}{$+0.0026$}
\newcommand{\SemiCellHighRevAGapD}{$+0.0094$$^{*}$}
\newcommand{\SemiCellHighRevBGapA}{$-0.0012$$^{*}$}
\newcommand{\SemiCellHighRevBGapB}{$+0.0022$$^{*}$}
\newcommand{\SemiCellHighRevBGapC}{$+0.0085$$^{*}$}
\newcommand{\SemiCellHighRevBGapD}{$+0.0139$$^{*}$}
\newcommand{\SemiCellHighRevCGapA}{$-0.0051$$^{*}$}
\newcommand{\SemiCellHighRevCGapB}{$-0.0013$}
\newcommand{\SemiCellHighRevCGapC}{$+0.0057$$^{*}$}
\newcommand{\SemiCellHighRevCGapD}{$+0.0120$$^{*}$}
\newcommand{\SemiZeroNSCells}{t=33, SD=0.05 (+0.0007)}

\IfFileExists{results/macros_robust.tex}{\newcommand{\RobSeeds}{20}
\newcommand{\RobAHPTenAllNeg}{20}
\newcommand{\RobAddLower}{18}
\newcommand{\RobAddSig}{5}
\newcommand{\RobAddMin}{$-$0.0007}
\newcommand{\RobAddMax}{0.0001}
\newcommand{\RobEWLower}{20}
\newcommand{\RobEWSig}{19}
\newcommand{\RobEWMin}{$-$0.0015}
\newcommand{\RobEWMax}{$-$0.0005}
}{}
% Semi-synthetic counts transcribed from the earlier semi-synthetic run (Table 10). \providecommand: a regenerated
% results/macros.tex (run_experiments.py --features ...) takes precedence automatically.
\providecommand{\SemiNegSigZero}{5}\providecommand{\SemiPosSigZero}{0}
\providecommand{\SemiNegSigHalf}{2}\providecommand{\SemiPosSigHalf}{2}
\providecommand{\SemiNegPointHalf}{3}\providecommand{\SemiPosPointHalf}{3}
\providecommand{\SemiPosPointOne}{6}\providecommand{\SemiPosSigOne}{5}
\providecommand{\SemiPosPointOneHalf}{6}\providecommand{\SemiPosSigOneHalf}{6}

\IfFileExists{results_empirical/macros_empirical.tex}{\newcommand{\EmpMonths}{32}
\newcommand{\EmpBestTen}{1.513}
\newcommand{\EmpBestFive}{2.026}
\newcommand{\EmpTenDiff}{$+0.065$}
\newcommand{\EmpTenLo}{$-0.075$}
\newcommand{\EmpTenHi}{$+0.187$}
\newcommand{\EmpFiveDiff}{$+0.146$}
\newcommand{\EmpFiveLo}{$+0.012$}
\newcommand{\EmpFiveHi}{$+0.289$}
\newcommand{\EmpSpearman}{$+0.007$}
\newcommand{\EmpFirstHalf}{1.716}
\newcommand{\EmpSecondHalf}{1.440}
\newcommand{\EmpWOcc}{0.53}
\newcommand{\EmpWForeign}{0.24}
\newcommand{\EmpWGuests}{0.13}
\newcommand{\EmpWReceipts}{0.08}
\newcommand{\EmpWVisitors}{0.02}
}{}

\newcommand{\RepoDOI}{10.5281/zenodo.XXXXXXX}   % <- paste the DOI of the corrected Zenodo release here
\newcommand{\PSDF}{\textsc{PSDF}}
\newcommand{\ie}{i.e.,\ }

\title{When Does Feedback-Driven Weighting Outperform Preference Elicitation?\\
Evidence from Simulation and Empirical Evaluation}
\author{Patcharaporn Samarnggoon\thanks{Independent Researcher, Chiang Mai, Thailand. E-mail: patcha535@gmail.com. ORCID: 0000-0003-3081-0605.}}
\date{}

\begin{document}
\maketitle

\begin{abstract}
\noindent\textbf{Purpose.} This study examines when a feedback-driven decision framework has an advantage over preference elicitation in sequential multi-criteria decision support. It introduces the Punch Score Decision Framework (PSDF), which estimates benefit-criterion weights from outcome feedback using rolling-window ridge regression and combines them with an explicit risk penalty, and evaluates it against static and periodically re-elicited Analytic Hierarchy Process (AHP) comparators.

\medskip\noindent\textbf{Design/methodology/approach.} A leakage-controlled simulation (100 replications, 30 alternatives, 60 periods) with a preference reversal is used. A $3\times4\times2$ factorial design varies reversal timing, the gap between stated and outcome-relevant preferences, and outcome noise. A semi-synthetic variant uses real 2024 Thai provincial tourism features. A separate empirical application ranks 77 Thai provinces monthly using Ministry of Tourism and Sports statistics from January 2023 to September 2026, with 32 decision months. Additional baselines include expanding-window, exponentially weighted and static ridge, CRITIC, uniform, signed ridge, momentum and random-selection rules. Simulation inference uses replication-level paired bootstrap intervals and Holm adjustment; empirical intervals use a moving-block bootstrap.

\medskip\noindent\textbf{Findings.} Under the baseline configuration, in which AHP receives noisy but unbiased current preference judgments, periodically re-elicited AHP ($k=10$) achieved lower expected regret than PSDF (mean paired difference, AHP minus PSDF, \DiffAHPTen, 95\% CI [\DiffAHPTenLo, \DiffAHPTenHi]). The ordering reverses when stated preferences differ systematically from the weights governing outcomes: at a preference gap of 1.5, PSDF had significantly lower regret in all six synthetic and all six semi-synthetic cells; at a gap of 1.0, its point estimate was lower in all twelve cells and statistically significant after Holm adjustment in \SynPosSigOne{} of 6 synthetic and \SemiPosSigOne{} of 6 semi-synthetic cells. PSDF regret was concentrated in the transition after the reversal. Adaptive weighting accounted for most of its advantage over non-adaptive weights, whereas the multiplicative and additive risk forms differed only marginally. An exponentially weighted ridge performed slightly better than the rolling-window estimator in repeated runs. In the empirical application, PSDF outperformed uniform and CRITIC weights but was not distinguishable from random selection at the top-10 level and was outperformed by a signed ridge forecast and a momentum rule.

\medskip\noindent\textbf{Originality/value.} The study reframes the comparison between feedback-driven weighting and AHP as a question of the conditions under which each information source is advantageous. It provides a regime map and separates adaptive weighting from risk aggregation. The factorial analysis is exploratory, the preference-gap magnitudes are not empirically calibrated, and the real-data application does not support a claim of predictive usefulness for PSDF as specified.

\medskip\noindent\textbf{Keywords:} Adaptive multi-criteria decision-making; decision support systems; feedback-driven weighting; Analytic Hierarchy Process; preference learning; sensitivity analysis.
\end{abstract}

\section{Introduction}
Multi-criteria decision-making (MCDM) provides a structured approach to evaluating alternatives when decisions involve multiple, potentially conflicting criteria. Such methods are widely used in decision support because they make the criteria, preferences, and aggregation rules underlying a recommendation more explicit.

A recurring challenge is that criterion importance may not remain constant over time. Decision contexts evolve, new information becomes available, and relevant preferences may change. Approaches that retain a fixed criterion-weight vector may therefore become less responsive to changes in the decision environment.

The Analytic Hierarchy Process (AHP) is a well-established approach for eliciting criterion priorities through pairwise comparisons \citep{saaty1980}. Its transparency and consistency diagnostics make it useful for structured decision problems. However, when preferences change, the decision maker must revisit the pairwise judgments or accept that the existing weights may no longer represent the current preference structure.

Research on dynamic MCDM and preference learning has addressed different aspects of this problem. Preference learning provides methods for estimating preference models from observed preference information \citep{bengs2021,hullermeier2024a,hullermeier2024b}, and recent work has investigated dynamic and hybrid criterion-weighting mechanisms under uncertainty \citep{ayan2023,lu2026,yang2025}.

This study proposes the Punch Score Decision Framework (PSDF), a feedback-driven decision-support framework that estimates benefit-criterion weights from historical outcome feedback and combines the resulting benefit score with an explicit risk penalty. It does not assume that outcome-derived weights are equivalent to weights elicited directly from decision makers; it treats them as distinct sources of information. The objective is not to establish a general claim of superiority over AHP but to identify the conditions under which each information source is advantageous.

The study addresses the following research questions.
\begin{description}
\item[RQ1.] How does rolling feedback-driven weighting behave when the underlying preference structure changes over time?
\item[RQ2.] How does adaptive weighting contribute to decision performance independently of the selected risk-aggregation function?
\item[RQ3.] How do static and periodically re-elicited AHP comparators behave under the same simulated decision environment?
\item[RQ4.] How sensitive is PSDF to incomplete outcome feedback and increasing observation noise?
\item[RQ5.] Under which combinations of reversal timing, stated--revealed preference gap, and noise does PSDF have an advantage over periodically re-elicited AHP?
\end{description}

The contributions are fourfold. (1) A formally specified feedback-driven decision framework that separates benefit weighting from risk aggregation. (2) A leakage-controlled sequential evaluation protocol with explicit temporal ordering and rolling-window rules. (3) Ablation, partial-feedback, noise-sensitivity, and additional-baseline experiments that examine the mechanisms contributing to observed decision performance. (4) A factorial regime map, in synthetic and semi-synthetic environments, that reports the conditions under which the ordering of PSDF and AHP changes, without interpreting differences between distinct feedback regimes as information-equivalent algorithmic superiority.

\section{Related Literature}
\subsection{Multi-criteria decision analysis and decision support}
MCDA is concerned not only with the mathematical aggregation of criteria but also with the broader process through which decision problems are structured, preferences are represented, and recommendations are produced. \citet{cinelli2020} provide a taxonomy of MCDA application processes and discuss the role of decision-support tools in method selection. AHP derives priorities from pairwise comparisons and provides a consistency diagnostic for the judgment matrix \citep{saaty1980}. It is particularly relevant when decision makers can express relative preferences but cannot directly specify a complete numerical weight vector. The present study uses AHP as a preference-elicitation comparator and considers both static and periodically re-elicited versions.

\subsection{Dynamic multi-criteria decision-making}
Dynamic MCDM extends conventional decision analysis to settings in which alternatives, criteria, or information change over time. Recent work has examined dynamic multi-criteria decision-making using learned information about alternatives \citep{yang2025}, and dynamic weighting has been investigated through hybrid subjective and objective approaches \citep{lu2026}. Reviews of weighting methods demonstrate the diversity of subjective, objective, and hybrid approaches \citep{ayan2023}. Adaptive weighting is therefore an established research direction, and this paper does not claim that feedback-driven weighting itself is a new concept.

\subsection{Preference learning and outcome-based feedback}
Preference learning estimates preference structures from observed preference information; its relationship to multiple-criteria decision aiding is examined in detail by \citet{hullermeier2024a,hullermeier2024b}. Online preference learning and dueling-bandit methods provide additional foundations for sequential preference estimation \citep{bengs2021}. However, the information used in preference learning may differ substantially from the information used in outcome prediction. This distinction is central to PSDF. The framework estimates criterion coefficients from observed outcome labels. These coefficients are predictive quantities and can be interpreted as decision-preference weights only under the assumption that the observed outcome adequately represents the decision objective. Predictive association alone does not establish normative preference importance.

\subsection{Positioning and research gap}
The proposed framework combines three design components: (i) rolling-window estimation of benefit-criterion weights from observed outcome feedback; (ii) explicit risk-adjusted score aggregation; and (iii) sequential evaluation against static and periodically re-elicited AHP. An observed performance difference may arise from adaptive weight estimation, risk aggregation, or their interaction, and a single comparison between complete systems cannot identify these contributions independently. The ablation design is intended to distinguish them, and the factorial design identifies conditions under which the comparison with AHP changes.

\section{Problem Formulation and Notation}
Consider a sequential decision problem with $m$ alternatives, $n$ benefit criteria, and $T$ decision periods. Let $Z_t\in[0,1]^{m\times n}$ denote the benefit-feature matrix at period $t$ with rows $Z_{it}$, and let $R_t=(R_{1t},\dots,R_{mt})^\top$, $R_{it}\in[0,1]$, denote the risk vector, where larger values represent greater risk. The true benefit-preference vector is $w^*_t=(w^*_{1t},\dots,w^*_{nt})^\top$ with $w^*_{jt}\ge0$ and $\sum_{j=1}^n w^*_{jt}=1$. The risk coefficient is $\rho>0$. The expected utility of alternative $i$ at period $t$ is
\begin{equation}\mu_{it}=Z_{it}w^*_t-\rho R_{it}.\end{equation}
The realized outcome and the observed feedback label are
\begin{align}U_{it}&=\mu_{it}+\varepsilon_{it},\\ Y_{it}&=U_{it}+\eta_{it},\end{align}
where $\varepsilon_{it}$ is outcome noise and $\eta_{it}$ is measurement or label noise, with separately specified distributions. Expected utility $\mu_{it}$ is the ground-truth decision criterion in the simulation; realized outcomes and labels are noisy observations and are not used as the benchmark for regret.

\section{The Punch Score Decision Framework}
\subsection{Risk-adjusted aggregation}
PSDF selects an alternative using the score
\begin{equation}S^{\mathrm{PSDF}}_{it}=\frac{Z_{it}\hat w_t}{1+\lambda R_{it}},\end{equation}
where $\hat w_t$ is the estimated benefit-weight vector and $\lambda\ge0$ controls the multiplicative risk penalty. The selected alternative is
\begin{equation}i^{\mathrm{PSDF}}_t=\arg\max_{i\in\{1,\dots,m\}}S^{\mathrm{PSDF}}_{it}.\end{equation}
The multiplicative formulation is not equivalent to an additive utility function. It is treated as a distinct aggregation rule and evaluated against the additive comparator
\begin{equation}S^{\mathrm{ADD}}_{it}=Z_{it}\hat w_t-\rho R_{it}.\end{equation}

\subsection{Rolling ridge estimation}
PSDF estimates benefit weights using completed-period feedback only. For each observed alternative-period record define $X_{is}=(Z_{i1s},\dots,Z_{ins},R_{is})^\top$. The rolling calendar window is $\mathcal W_t=\{(i,s):\max(0,t-L)\le s<t,\ Y_{is}\text{ is observed}\}$. The ridge regression with an unpenalized intercept solves
\begin{equation}(\hat\beta_{0,t},\hat\beta_t)=\arg\min_{\beta_0,\beta}\Big\{\sum_{(i,s)\in\mathcal W_t}\big(Y_{is}-\beta_0-X_{is}^\top\beta\big)^2+\alpha\|\beta\|_2^2\Big\}.\end{equation}
The benefit coefficients are the first $n$ components of $\hat\beta_t$; the risk coefficient is retained as a predictor but is not converted into a benefit weight. The non-negative normalized benefit weights are
\begin{equation}\hat w_{jt}=\frac{\max(0,\hat\beta_{jt})}{\sum_{\ell=1}^n\max(0,\hat\beta_{\ell t})}.\end{equation}
The estimator is fitted only when at least $N_{\min}=12$ observed records are available in the window. If fewer records are available, or the normalization denominator is zero or numerically unstable (below $10^{-12}$), uniform weights are used:
\begin{equation}\hat w_{jt}=\frac1n.\end{equation}
This fallback is specified in advance and is not selected according to test-period performance.

\subsection{Rolling-window specification}
The primary window is $L=15$ completed calendar periods. Under partial feedback, only observed labels are included and missing labels are not imputed. At $t=0$ no labels exist, so $\hat w_0=(1/n,\dots,1/n)^\top$. Feedback generated at period $t$ can first influence the decision at period $t+1$.

\subsection{Fixed-weight comparator initialization}
The fixed-weight PSDF comparator (also called static ridge) uses the first eligible ridge estimate produced after period-zero feedback (the estimate used at $t=1$) and retains it unchanged thereafter; the period-zero decision uses uniform weights. If the minimum-observation condition is not met, the uniform fallback is retained. This comparator isolates the effect of subsequent weight adaptation from the multiplicative aggregation rule and is not an oracle with access to the true initial preference vector.

\subsection{Interpretation of estimated weights}
The regression coefficients are outcome-predictive coefficients, and their conversion into non-negative normalized weights imposes a decision representation on them. This is meaningful only if the observed outcome is a defensible proxy for the decision objective. If outcomes are affected by omitted variables, selection effects, or feedback loops, the estimated weights may not recover normative preferences. Recovery of a known preference vector in simulation does not establish that real-world outcome feedback identifies human preferences.

\section{AHP Comparator}\label{sec:ahp}
\subsection{Criterion representation}
The AHP comparator uses the $n$ benefit criteria and one safety criterion $Q_{i,n+1,t}=1-R_{it}$ (six criteria in the synthetic environment with $n=5$; four in the semi-synthetic environment with $n=3$). The target AHP priority vector is
\begin{equation}v^*_t=\frac{(w^*_{1t},\dots,w^*_{nt},\rho)^\top}{1+\rho},\end{equation}
and the weighted AHP rating is
\begin{equation}S^{\mathrm{AHP}}_{it}=\sum_{j=1}^n v_{jt}Z_{ijt}+v_{n+1,t}(1-R_{it}).\end{equation}
When $v_t=v^*_t$, this score is a positive affine transformation of expected utility, so the comparator explicitly incorporates risk. This construction does not make the feedback channels equivalent.

\subsection{Pairwise comparison matrix}
At each scheduled elicitation round $\tau$, the simulated decision maker provides pairwise judgments for all unordered pairs of criteria. For $j<\ell$,
\begin{equation}a_{j\ell,\tau}=\operatorname{clip}\Big[\frac{v_{j,\tau}}{v_{\ell,\tau}}\exp(\epsilon_{j\ell,\tau}),\ \tfrac19,\ 9\Big],\qquad \epsilon_{j\ell,\tau}\sim N(0,0.05^2),\end{equation}
with reciprocal entries $a_{\ell j,\tau}=1/a_{j\ell,\tau}$ and unit diagonal. Here $v_\tau$ is the stated priority vector at round $\tau$: in the baseline it equals the true current priority vector $v^*_\tau$, whereas in the preference-gap experiments of Section~\ref{sec:factorial} the benefit weights are distorted before judgments are generated. The perturbation is generated once per unordered pair. The judgments therefore represent a simulated elicitation channel with access to current preference information at scheduled rounds. This is an explicit modelling assumption, not an assertion that real decision makers can reveal preferences without error or cost.

\subsection{Priority estimation and consistency diagnostics}
Priorities are obtained from the principal right eigenvector,
\begin{equation}A_\tau v_\tau=\lambda_{\max,\tau}v_\tau,\qquad\sum_{j=1}^{n+1}v_{j,\tau}=1.\end{equation}
The consistency index and ratio are
\begin{equation}\mathrm{CI}_\tau=\max\Big(0,\frac{\lambda_{\max,\tau}-n_A}{n_A-1}\Big),\qquad \mathrm{CR}_\tau=\frac{\mathrm{CI}_\tau}{\mathrm{RI}_{n_A}},\end{equation}
where $n_A=n+1$, $\mathrm{RI}_6=1.24$ for the six-criterion matrix and $\mathrm{RI}_4=0.90$ for the four-criterion matrix. The study reports the mean, standard deviation, median and interquartile range of CR, the proportion of matrices with $\mathrm{CR}\le0.10$, and the proportion of pairwise judgments affected by the $1/9$ and $9$ clipping limits. The $\mathrm{CR}\le0.10$ convention is a diagnostic threshold, not proof that judgments are correct or unbiased.

\subsection{Periodic re-elicitation}
AHP weights are updated only at scheduled rounds,
\begin{equation}\tau_k(t)=k\lfloor t/k\rfloor,\qquad v_t=v_{\tau_k(t)},\end{equation}
and the most recent priority vector is held between rounds. The primary schedule is $k=10$, with sensitivity settings $k=5$ and $k=20$. Static AHP is a single elicitation at $t=0$ with no later updates. Because the baseline reversal occurs at $t=30$, rounds for $k=5$ and $k=10$ coincide with the reversal, whereas $k=20$ re-elicits at $t=40$.

\subsection{AHP-weighted PSDF}
A weight-source ablation uses the benefit priorities from the initial static AHP elicitation,
\begin{equation}\hat w^{\mathrm{AHP},0}_j=\frac{v_{j,0}}{\sum_{\ell=1}^n v_{\ell,0}},\qquad S^{\mathrm{AHP\text{-}PSDF}}_{it}=\frac{Z_{it}\hat w^{\mathrm{AHP},0}}{1+\lambda R_{it}},\end{equation}
inside the multiplicative PSDF score. This configuration is not periodically re-elicited. Comparing it with static AHP examines two aggregation rules using the same initial AHP preference source.

\section{Temporal Integrity and Reproducibility}
\subsection{Information timing}
The simulation follows a strict sequence: (1) generate the current benefit and risk features; (2) estimate PSDF weights using completed historical feedback only; (3) retrieve the latest available AHP priority vector under the elicitation schedule; (4) calculate method-specific scores and select alternatives; (5) calculate expected-utility regret; (6) generate realized outcomes and observed labels; and (7) add available feedback to the historical set for subsequent periods. No method may use current-period feedback to select an alternative in that period.

\subsection{Temporal-integrity protocol}
The implementation is checked by an automated test suite (supplementary material): no future-label access (corrupting labels from period $t_0$ onward leaves all decisions up to $t_0$ unchanged, while later decisions change); no future-judgment access; AHP weight freezing between rounds; rejection of out-of-order period updates; immutability of historical snapshots; feedback timing (current-period labels enter only after the decision); the minimum-observation fallback; the partial-feedback counts; and reproducibility under the stored seed. Passing software tests supports implementation integrity but does not establish external validity.

\subsection{Algorithm}
\begin{algorithm}[t]
\caption{Sequential PSDF and AHP evaluation protocol}
\begin{algorithmic}[1]
\Require Alternatives $m$, criteria $n$, periods $T$, window $L$, minimum observations $N_{\min}$, ridge penalty $\alpha$, risk parameter $\lambda$, AHP interval $k$
\State Initialize historical feedback set $H_0=\emptyset$ and $\hat w_0=\mathbf 1/n$
\State Elicit initial AHP pairwise matrix and calculate $v_0$
\For{$t=0,\dots,T-1$}
  \State Generate current $Z_t$ and $R_t$
  \State Retain observed feedback from periods $\max(0,t-L),\dots,t-1$
  \If{number of observed records $\ge N_{\min}$}
    \State Fit ridge regression with unpenalized intercept
    \State Normalize non-negative benefit slopes; fall back to uniform if normalization fails
  \Else
    \State Set PSDF weights to uniform weights
  \EndIf
  \If{$t>0$ and $t\bmod k=0$}
    \State Elicit current AHP judgments; compute priorities, $\mathrm{CI}_t$ and $\mathrm{CR}_t$
  \EndIf
  \State Calculate PSDF and comparator scores; select one alternative per method
  \State Calculate expected-utility regret
  \State Generate realized outcomes and labels; add eligible feedback for period $t+1$
\EndFor
\State \Return Decisions, weights, regret, and diagnostics
\end{algorithmic}
\end{algorithm}

\section{Simulation Design}
\subsection{Base simulation configuration}
The base configuration is shown in Table~\ref{tab:config}. The true preference vector is
\begin{equation}w^*_t=\begin{cases}(0.70,0.10,0.08,0.07,0.05)^\top,&t<30,\\(0.05,0.07,0.08,0.10,0.70)^\top,&t\ge30.\end{cases}\end{equation}
Benefit features and risk are independent Beta(2,2) draws. Replication-specific random streams are derived from the master seed by a seed-sequence spawn and are separated into environment generation, outcome/label noise, feedback sampling, AHP judgment noise, and the preference-gap direction. Within a replication, all methods and all experimental cells are evaluated on the same generated environment.

\begin{table}[t]
\centering\small
\caption{Base simulation configuration.}\label{tab:config}
\begin{tabular}{ll}
\toprule
Parameter & Specification\\
\midrule
Number of alternatives & 30\\
Number of benefit criteria & 5\\
Number of periods & 60\\
Number of replications & 100\\
Preference reversal & Period 30\\
Benefit-feature distribution & Beta(2,2)\\
Risk distribution & Beta(2,2)\\
Initial benefit weights & (0.70, 0.10, 0.08, 0.07, 0.05)\\
Post-reversal weights & (0.05, 0.07, 0.08, 0.10, 0.70)\\
Risk coefficient $\rho$ & 0.5\\
PSDF risk parameter $\lambda$ & 1.0\\
Ridge penalty $\alpha$ & 0.1\\
Rolling-window length $L$ & 15 periods\\
Minimum regression observations & 12\\
Outcome-noise SD & 0.05\\
Label-noise SD & 0.05\\
AHP judgment log-noise SD & 0.05\\
AHP primary interval & 10 periods\\
AHP sensitivity intervals & 5 and 20 periods\\
Master random seed & 6621\\
\bottomrule
\end{tabular}
\end{table}

\subsection{Decision methods and comparators}
The primary methods are: (1) adaptive PSDF (rolling ridge weights, multiplicative risk); (2) adaptive additive PSDF (same weights, additive risk); (3) fixed-weight PSDF (weights frozen after period-zero feedback, multiplicative risk); (4) static AHP (risk-inclusive, elicited at $t=0$); (5) periodic AHP (risk-inclusive, $k=5,10,20$); and (6) AHP-weighted PSDF (initial static AHP benefit weights in the multiplicative score), which is an aggregation ablation and not an independently updated AHP method. Four additional comparators are used in the sensitivity analysis, all with the multiplicative score: (7) an expanding-window ridge using all completed periods; (8) an exponentially weighted ridge with discount factor 0.9 per period (record weight $0.9^{t-1-s}$) and the same $\alpha$ and $N_{\min}$; (9) CRITIC weights computed from the current feature matrix (min--max normalized columns) without outcome feedback; and (10) uniform weights. Two diagnostic references that are not competing methods are also computed: the multiplicative score evaluated with the true weights, and the additive score evaluated with the true weights.

\subsection{Ablation design}
\begin{table}[t]
\centering\small
\caption{Ablation experiments and their analytical purpose.}
\begin{tabular}{p{3.0cm}p{6.3cm}p{4.6cm}}
\toprule
Ablation & Comparison & Purpose\\
\midrule
Weight adaptation & Adaptive versus fixed benefit weights under the same multiplicative score & Contribution of adaptive weighting\\
Risk aggregation & Multiplicative versus additive aggregation with the same estimated weights & Contribution of aggregation form\\
Weight-source control & Initial AHP-derived versus feedback-derived benefit weights in the PSDF score & Separate weight source from aggregation rule\\
AHP update frequency & Static AHP versus periodic AHP at $k=5,10,20$ & Sensitivity to elicitation frequency\\
Risk-inclusive AHP & AHP with the explicit safety criterion $1-R$ & Avoid confounding through omission of risk\\
\bottomrule
\end{tabular}
\end{table}
The ablations are controlled comparisons within the specified environment and do not identify causal effects in external settings.

\subsection{Partial-feedback experiment}
The base configuration assumes full feedback. In the partial-feedback experiment, let $q\in\{0.25,0.50,0.75,1.00\}$ be the target proportion of alternatives whose labels are observed after each period. Observed alternatives are sampled uniformly without replacement within each period, and unobserved labels are not imputed. Because the number of alternatives is discrete, the number observed per period is $\operatorname{round}(qm)$ with ties rounded to the even integer, \ie 8, 15, 22 and 30 of 30 alternatives (achieved proportions 26.7\%, 50.0\%, 73.3\% and 100\%). The condition $q=1$ is the reference.

\subsection{Factorial regime map}\label{sec:factorial}
To identify when the ordering of PSDF and AHP changes, three factors are crossed: (i) the reversal period $\in\{30,33,37\}$, which determines alignment with the AHP schedule ($k=10$ re-elicits at multiples of ten); (ii) a persistent preference gap $g\in\{0,0.5,1.0,1.5\}$, in which the stated benefit weights used to generate AHP judgments equal $w^*_t\odot e^{gz}$, renormalized, with a fixed $z\sim N(0,I)$ per replication (the safety weight is not distorted); and (iii) outcome and label noise SD $\in\{0.05,0.20\}$. The gap represents a systematic difference between stated preferences and the weights that govern outcomes; regret is always measured against the true expected utility. The factorial was specified after the baseline result was known and is therefore exploratory.

\subsection{Semi-synthetic environment}
In the semi-synthetic variant, each period draws 30 of the 77 Thai provinces without replacement, and the benefit features are real 2024 provincial values from the revised Ministry of Tourism and Sports workbook \citep{mots2025}: the logarithm of recorded visits, the logarithm of receipts per recorded visit, and the foreign share of recorded visits, each min--max scaled to $[0,1]$. The logarithm of receipts is omitted because it is an exact linear combination of the first two. There are therefore $n=3$ benefit criteria, with true weights $(0.60,0.25,0.15)$ reversing to $(0.15,0.25,0.60)$. Risk, preferences and outcomes remain simulated, so this variant tests sensitivity to a realistic, heavy-tailed feature distribution and not empirical validity. Absolute regret values are not comparable across the synthetic and semi-synthetic environments because the criterion counts differ.

\subsection{Noise and Monte Carlo stability experiments}
In the noise-sensitivity experiment, outcome and label noise are varied jointly (the same SD, $\in\{0,0.05,0.10,0.20\}$) while AHP judgment noise is held at 0.05; AHP does not use outcome feedback and is therefore unaffected by this manipulation. To assess which conclusions depend on the particular random draws, the baseline experiment is repeated under 20 independent master seeds (1001--1020) and the paired contrasts are summarized across runs (Section~\ref{sec:stability}).

\section{Evaluation Metrics}
\subsection{Expected-utility regret}
Regret is measured against the oracle alternative $i^{\mathrm{oracle}}_t=\arg\max_i\mu_{it}$:
\begin{equation}\mathrm{Regret}_t=\max_i\mu_{it}-\mu_{i_t,t},\qquad \mathrm{Regret}_r=\frac1T\sum_{t=1}^T\mathrm{Regret}_{rt}.\end{equation}
Expected-utility regret is used instead of realized-outcome regret because the latter includes random noise that is not part of the decision rule's objective.

\subsection{Weight recovery}
Weight recovery is measured by the $L_1$ distance between estimated and true benefit weights,
\begin{equation}E^w_{rt}=\sum_{j=1}^n\big|\hat w_{rjt}-w^*_{rjt}\big|.\end{equation}
AHP benefit weights are normalized over the benefit criteria before comparison; the safety weight is reported separately. The pre-reversal period is $t=0,\dots,29$ and the post-reversal period is $t=30,\dots,59$; the uniform-weight period $t=0$ is retained in the full-horizon summary.

\section{Statistical Analysis}
The independent inferential unit is the simulation replication. For methods A and B define
\begin{equation}d_r=\mathrm{Regret}_{A,r}-\mathrm{Regret}_{B,r},\end{equation}
so that a positive value indicates lower regret for method B. The analysis reports the mean paired difference, a 95\% paired percentile bootstrap confidence interval (10{,}000 resamples of complete replications in the baseline analysis, 5{,}000 in the factorial cells), a two-sided paired $t$-test, and Holm-adjusted $p$-values.

The prespecified family for the baseline analysis contains seven paired comparisons: (1) periodic AHP ($k=10$) versus adaptive PSDF; (2) static AHP versus adaptive PSDF; (3) fixed-weight PSDF versus adaptive PSDF; (4) additive versus multiplicative PSDF; (5) periodic AHP ($k=5$) versus adaptive PSDF; (6) periodic AHP ($k=20$) versus adaptive PSDF; and (7) AHP-weighted PSDF versus adaptive PSDF. The sign of each reported difference follows the order shown in the corresponding results table; reversing the order changes the sign but not the two-sided $p$-value.

For the factorial regime map, the effect is $d=\mathrm{Regret}_{\mathrm{AHP}(k=10)}-\mathrm{Regret}_{\mathrm{PSDF}}$ in each cell, and Holm adjustment is applied across the 24 cells of each environment; positive values favor PSDF. The null hypothesis is $H_0:E[d_r]=0$ against a two-sided alternative. Direction and magnitude are reported regardless of significance, and no method is called superior solely because its point estimate is smaller. Comparisons with the additional baselines and the partial-feedback and noise experiments are exploratory and not multiplicity-adjusted.

\section{Results}\label{sec:results}
The simulation results in this section were produced by the code in the supplementary material under the configuration in Table~\ref{tab:config} (run identifier \texttt{\RunID}). The empirical application is reported separately in Section~\ref{sec:empirical}.

\subsection{Primary performance comparison}
Periodic AHP ($k=10$) achieved the lowest expected regret, followed by adaptive PSDF and then static AHP (Table~\ref{tab:primary}). AHP regret is close to zero because its simulated judgments are generated from the current true preference vector with small noise (Section~\ref{sec:ahp}).

\begin{table}[t]
\centering\small
\caption{Primary expected-regret comparison (100 replications; bootstrap 95\% CI of the mean).}\label{tab:primary}
\begin{tabular}{lrrrr}
\toprule
Method & Mean regret & 95\% CI lower & 95\% CI upper & Replication SD \\
\midrule
Adaptive PSDF (rolling ridge) & 0.0202 & 0.0189 & 0.0214 & 0.0064 \\
Periodic AHP ($k=10$) & 0.0002 & 0.0001 & 0.0002 & 0.0002 \\
Static AHP & 0.0945 & 0.0917 & 0.0974 & 0.0146 \\
\bottomrule
\end{tabular}

\end{table}

By phase, mean regret before the reversal was \PreAdaptive{} for adaptive PSDF and \PreAHPTen{} for periodic AHP ($k=10$). After the reversal, adaptive PSDF averaged \PostAdaptive{} (\TransAdaptive{} over $t=30$--44 and \LateAdaptive{} over $t=45$--59), whereas periodic AHP remained at \PostAHPTen{}. The PSDF deficit is concentrated in the transition while the rolling window still contains pre-reversal observations.

\subsection{Paired statistical comparisons}
Periodic AHP with $k=5$ and $k=10$ achieved lower regret than adaptive PSDF, whereas periodic AHP with $k=20$ had higher regret because its first post-reversal elicitation occurs at $t=40$. Adaptive PSDF achieved lower regret than static AHP, fixed-weight PSDF and AHP-weighted PSDF. The additive--multiplicative difference was \AddDiff{} (Holm $p=\AddHolm$), small relative to the level of regret (\RegAdaptive); its stability across random seeds is examined in Section~\ref{sec:stability}.

\begin{table}[t]
\centering\small
\caption{Paired inference for prespecified comparisons (mean difference of replication-level mean regret, first minus second method). A negative difference indicates lower mean regret for the first-named method; a positive difference indicates lower regret for adaptive PSDF.}\label{tab:paired}
\begin{tabular}{lrcrr}
\toprule
Comparison & Mean diff. & 95\% CI & Paired $p$ & Holm $p$ \\
\midrule
Periodic AHP ($k=10$) $-$ adaptive PSDF & $-0.0200$ & [$-$0.0213, $-$0.0188] & $6.0\times10^{-53}$ & $2.4\times10^{-52}$ \\
Static AHP $-$ adaptive PSDF & $+0.0743$ & [0.0714, 0.0773] & $1.3\times10^{-71}$ & $9.0\times10^{-71}$ \\
Fixed-weight PSDF $-$ adaptive PSDF & $+0.0794$ & [0.0763, 0.0825] & $2.5\times10^{-71}$ & $1.5\times10^{-70}$ \\
Additive PSDF $-$ multiplicative PSDF & $-0.0005$ & [$-$0.0009, $-$0.0001] & 0.025 & 0.025 \\
Periodic AHP ($k=5$) $-$ adaptive PSDF & $-0.0200$ & [$-$0.0213, $-$0.0188] & $6.2\times10^{-53}$ & $2.4\times10^{-52}$ \\
Periodic AHP ($k=20$) $-$ adaptive PSDF & $+0.0121$ & [0.0107, 0.0135] & $3.0\times10^{-30}$ & $6.1\times10^{-30}$ \\
AHP-weighted PSDF $-$ adaptive PSDF & $+0.0741$ & [0.0711, 0.0771] & $1.7\times10^{-70}$ & $8.3\times10^{-70}$ \\
\bottomrule
\end{tabular}

\end{table}

\subsection{Weight-recovery analysis}
\begin{table}[t]
\centering\small
\caption{Weight-recovery results (mean $L_1$ error of normalized benefit weights).}\label{tab:weights}
\begin{tabular}{lrrr}
\toprule
Method & Pre-reversal & Post-reversal & Overall \\
\midrule
Adaptive PSDF & 0.113 & 0.406 & 0.260 \\
Periodic AHP ($k=10$) & 0.046 & 0.047 & 0.047 \\
Periodic AHP ($k=5$) & 0.046 & 0.046 & 0.046 \\
Periodic AHP ($k=20$) & 0.047 & 0.475 & 0.261 \\
Static AHP & 0.048 & 1.333 & 0.690 \\
Fixed-weight PSDF & 0.268 & 1.431 & 0.850 \\
\bottomrule
\end{tabular}

\end{table}
Adaptive PSDF starts with an $L_1$ error of \LOneZero{} at $t=0$ (uniform weights against a concentrated true vector). Its error is low before the reversal and is driven by the post-reversal transition afterwards. AHP errors reflect judgment noise and clipping of extreme ratios.

\subsection{Ablation analysis}
\begin{table}[t]
\centering\small
\caption{Ablation results (paired difference relative to adaptive weights with multiplicative risk).}\label{tab:ablation}
\begin{tabular}{lccr}
\toprule
Configuration & Mean regret & 95\% CI & Paired difference \\
\midrule
Adaptive weights, multiplicative risk & 0.0202 & [0.0190, 0.0214] & -- \\
Fixed weights, multiplicative risk & 0.0996 & [0.0964, 0.1027] & $+0.0794$ \\
Adaptive weights, additive risk & 0.0197 & [0.0184, 0.0211] & $-0.0005$ \\
Initial AHP weights, multiplicative PSDF & 0.0943 & [0.0914, 0.0973] & $+0.0741$ \\
\bottomrule
\end{tabular}

\end{table}
Holding aggregation constant, adaptive weights yield markedly lower regret than weights frozen after period zero or taken from the initial AHP elicitation. Holding weights constant, the aggregation form has a very small effect (Table~\ref{tab:ablation}); a statistically detectable but practically negligible difference is not evidence of equivalence, and the contrast lies near the boundary of detectability across seeds (Section~\ref{sec:stability}). As a diagnostic, the multiplicative score evaluated with the \emph{true} weights has mean regret \OracleMultRegret{} (additive: \OracleAddRegret{}), which is the regret attributable to the multiplicative form alone.

\subsection{Partial-feedback analysis}
\begin{table}[t]
\centering\small
\caption{Partial-feedback results for adaptive PSDF.}\label{tab:partial}
\begin{tabular}{lrrr}
\toprule
Target rate (observed per period) & Mean expected regret & 95\% CI lower & 95\% CI upper \\
\midrule
25\% (8 of 30; achieved 26.7\%) & 0.0227 & 0.0214 & 0.0240 \\
50\% (15 of 30; achieved 50.0\%) & 0.0209 & 0.0196 & 0.0222 \\
75\% (22 of 30; achieved 73.3\%) & 0.0205 & 0.0193 & 0.0217 \\
100\% (30 of 30; achieved 100\%) & 0.0202 & 0.0190 & 0.0215 \\
\bottomrule
\end{tabular}

\end{table}
Performance degraded modestly as feedback was reduced, with the largest loss at the lowest feedback rate. The nominal feedback rates are targets; the achieved rates reflect the integer number of alternatives observed in each period. At the 25\% rate the minimum-observation rule occasionally forces the uniform fallback early in the horizon.

\subsection{Noise sensitivity}
Outcome and label noise are varied jointly, holding AHP judgment noise at its base value (Table~\ref{tab:noise}).
\begin{table}[t]
\centering\small
\caption{Noise sensitivity: mean expected regret.}\label{tab:noise}
\begin{tabular}{lrrrc}
\toprule
Noise SD & Adaptive PSDF & Periodic AHP ($k=10$) & Difference (AHP $-$ PSDF) & 95\% CI \\
\midrule
0.00 & 0.0198 & 0.0002 & $-0.0196$ & [$-$0.0208, $-$0.0184] \\
0.05 & 0.0202 & 0.0002 & $-0.0200$ & [$-$0.0213, $-$0.0188] \\
0.10 & 0.0214 & 0.0002 & $-0.0212$ & [$-$0.0225, $-$0.0199] \\
0.20 & 0.0238 & 0.0002 & $-0.0237$ & [$-$0.0252, $-$0.0222] \\
\bottomrule
\end{tabular}

\end{table}
PSDF regret rose gradually with noise, from \NoisePSDFZero{} to \NoisePSDFHigh{}, and the ordering relative to periodic AHP was unchanged at the baseline gap and reversal timing.

\subsection{AHP consistency diagnostics}
For the $k=10$ schedule (\CRMatrices{} matrices: 6 rounds $\times$ 100 replications), mean CR was \CRMean{} (SD \CRSD{}; median \CRMedian{}, IQR \CRQOne--\CRQThree{}), and the proportion of matrices satisfying $\mathrm{CR}\le0.10$ was \CRProp{}. Clipping at $1/9$ and $9$ affected \ClipN{} of \ClipTotal{} pairwise judgments (\ClipPct\%), mainly ratios involving the dominant criterion and the safety criterion: with the baseline weights, three of the 15 target ratios exceed the clipping limits before any noise is added. These diagnostics describe internal consistency only and do not establish that the underlying preference vector is normatively correct.

\subsection{Sensitivity analysis: when does PSDF have an advantage?}
The primary comparison favors periodic AHP under the baseline configuration. The factorial design in Section~\ref{sec:factorial} asks under which conditions the ordering changes. It is exploratory; all 24 cells are shown per environment, including those in which AHP is superior. The effect is $d=\mathrm{Regret}_{\mathrm{AHP}(k=10)}-\mathrm{Regret}_{\mathrm{PSDF}}$, so positive values favor PSDF.

\begin{table}[t]
\centering\small
\caption{Regime map, synthetic environment: $d=\mathrm{Regret}_{\mathrm{AHP}}-\mathrm{Regret}_{\mathrm{PSDF}}$ ($k=10$; positive favors PSDF). An asterisk marks cells with Holm-adjusted $p<0.05$ across the 24 cells.}\label{tab:regime-syn}
\begin{tabular}{llrrrr}
\toprule
Noise SD & Reversal & $g=0$ & $g=0.5$ & $g=1.0$ & $g=1.5$ \\
\midrule
0.05 & $t=30$ & $-0.0200$$^{*}$ & $-0.0141$$^{*}$ & $+0.0071$ & $+0.0314$$^{*}$ \\
0.05 & $t=33$ & $+0.0022$$^{*}$ & $+0.0065$$^{*}$ & $+0.0236$$^{*}$ & $+0.0447$$^{*}$ \\
0.05 & $t=37$ & $-0.0108$$^{*}$ & $-0.0055$$^{*}$ & $+0.0139$$^{*}$ & $+0.0368$$^{*}$ \\
0.20 & $t=30$ & $-0.0237$$^{*}$ & $-0.0177$$^{*}$ & $+0.0034$ & $+0.0278$$^{*}$ \\
0.20 & $t=33$ & $-0.0016$ & $+0.0027$ & $+0.0198$$^{*}$ & $+0.0409$$^{*}$ \\
0.20 & $t=37$ & $-0.0145$$^{*}$ & $-0.0092$$^{*}$ & $+0.0102$$^{*}$ & $+0.0332$$^{*}$ \\
\bottomrule
\end{tabular}

\end{table}

\begin{table}[t]
\centering\small
\caption{Regime map, semi-synthetic environment (real provincial features). Same layout as Table~\ref{tab:regime-syn}.}\label{tab:regime-semi}
\IfFileExists{results/tex/table10.tex}{\begin{tabular}{llrrrr}
\toprule
Noise SD & Reversal & $g=0$ & $g=0.5$ & $g=1.0$ & $g=1.5$ \\
\midrule
0.05 & $t=30$ & $-0.0060$$^{*}$ & $-0.0024$$^{*}$ & $+0.0047$$^{*}$ & $+0.0115$$^{*}$ \\
0.05 & $t=33$ & $+0.0007$ & $+0.0041$$^{*}$ & $+0.0104$$^{*}$ & $+0.0158$$^{*}$ \\
0.05 & $t=37$ & $-0.0028$$^{*}$ & $+0.0011$ & $+0.0080$$^{*}$ & $+0.0143$$^{*}$ \\
0.20 & $t=30$ & $-0.0080$$^{*}$ & $-0.0044$$^{*}$ & $+0.0026$ & $+0.0094$$^{*}$ \\
0.20 & $t=33$ & $-0.0012$$^{*}$ & $+0.0022$$^{*}$ & $+0.0085$$^{*}$ & $+0.0139$$^{*}$ \\
0.20 & $t=37$ & $-0.0051$$^{*}$ & $-0.0013$ & $+0.0057$$^{*}$ & $+0.0120$$^{*}$ \\
\bottomrule
\end{tabular}
}{\begin{tabular}{llrrrr}
\toprule
Noise SD & Reversal & $g=0$ & $g=0.5$ & $g=1.0$ & $g=1.5$ \\
\midrule
0.05 & $t=30$ & $-0.0059^{*}$ & $-0.0019^{*}$ & $+0.0052^{*}$ & $+0.0120^{*}$ \\
0.05 & $t=33$ & $+0.0007$ & $+0.0043^{*}$ & $+0.0108^{*}$ & $+0.0164^{*}$ \\
0.05 & $t=37$ & $-0.0029^{*}$ & $+0.0011$ & $+0.0081^{*}$ & $+0.0143^{*}$ \\
0.20 & $t=30$ & $-0.0078^{*}$ & $-0.0037^{*}$ & $+0.0033$ & $+0.0101^{*}$ \\
0.20 & $t=33$ & $-0.0013^{*}$ & $+0.0023^{*}$ & $+0.0087^{*}$ & $+0.0144^{*}$ \\
0.20 & $t=37$ & $-0.0049^{*}$ & $-0.0008$ & $+0.0061^{*}$ & $+0.0123^{*}$ \\
\bottomrule
\end{tabular}
}
\end{table}

\begin{table}[t]
\centering\small
\caption{Mean expected regret for all methods in four representative cells (synthetic; noise SD 0.05 except the last column).}\label{tab:allmethods}
\begin{tabular}{lrrrr}
\toprule
Method & $t=30$, $g=0$ & $t=37$, $g=0$ & $t=30$, $g=1$ & $t=37$, $g=1$, SD 0.20 \\
\midrule
Adaptive PSDF (rolling ridge) & 0.0202 & 0.0208 & 0.0202 & 0.0244 \\
Adaptive additive PSDF & 0.0197 & 0.0205 & 0.0197 & 0.0241 \\
Exponentially weighted ridge & 0.0197 & 0.0203 & 0.0197 & 0.0230 \\
Expanding-window ridge & 0.0553 & 0.0518 & 0.0553 & 0.0533 \\
Static ridge (after period 0) & 0.0996 & 0.0781 & 0.0996 & 0.0862 \\
CRITIC weights & 0.0839 & 0.0846 & 0.0839 & 0.0846 \\
Uniform weights & 0.0825 & 0.0832 & 0.0825 & 0.0832 \\
Static AHP & 0.0945 & 0.0724 & 0.1000 & 0.0823 \\
Periodic AHP ($k=10$) & 0.0002 & 0.0100 & 0.0273 & 0.0347 \\
Periodic AHP ($k=20$) & 0.0323 & 0.0100 & 0.0524 & 0.0348 \\
AHP-weighted PSDF & 0.0943 & 0.0723 & 0.0982 & 0.0812 \\
\bottomrule
\end{tabular}

\end{table}

\paragraph{Findings.} Three patterns are visible. First, when AHP judgments are unbiased ($g=0$), periodic AHP has significantly lower regret in \SynNegSigZero{} of 6 synthetic and \SemiNegSigZero{} of 6 semi-synthetic cells. The two cells with the reversal at $t=33$ lie within a few thousandths of zero: in the synthetic environment the difference is \SynCellLowRevBGapA{} at low noise and \SynCellHighRevBGapA{} at high noise ($^{*}$ marks Holm-significance), so the sign and significance of these near-zero cells should not be over-interpreted. Second, as the preference gap grows the ordering reverses. At $g=1.0$ the point estimate favors PSDF in all six synthetic and all six semi-synthetic cells, with Holm-adjusted significance in \SynPosSigOne{} of 6 synthetic and \SemiPosSigOne{} of 6 semi-synthetic cells; at $g=1.5$ PSDF has significantly lower regret in all twelve cells. At $g=0.5$ the picture is mixed: AHP is significantly better in \SynNegSigHalf{} of 6 synthetic cells but in only \SemiNegSigHalf{} of 6 semi-synthetic cells. Third, misalignment between the reversal and the elicitation schedule matters: at $g=0$ and low noise the AHP advantage shrinks from \AHPAdvEarly{} ($t=30$) to \AHPShrinkLate{} ($t=37$) and is close to zero at $t=33$.

Among outcome-based methods, an exponentially weighted ridge had a marginally lower regret than the rolling-window ridge in the baseline run (EW minus rolling: \EWDiff, 95\% CI [\EWLo, \EWHi]); the repeated-seed analysis below shows that this small advantage is consistent. The rolling window is therefore one of several reasonable ways to discount old feedback and not a uniquely favorable design. Expanding-window, static, CRITIC and uniform weights were all substantially worse than adaptive PSDF in the representative cells shown (Table~\ref{tab:allmethods}), supporting the interpretation that discounting of old outcomes drives PSDF's performance relative to non-adaptive alternatives. At $t=37$ with $g=0$, periodic AHP with $k=10$ and $k=20$ coincide because both first update at $t=40$.

\paragraph{Caveats.} The preference-gap factor is a modelling construct. Its magnitude ($g=1.0$ corresponds to stated weights distorted by a typical factor of $e^{\pm1}$ per criterion) has not been calibrated against observed discrepancies between stated and revealed preferences, and at $g=0.5$ AHP remains superior in most synthetic cells. The factorial was specified after the baseline result was known.

\subsection{Monte Carlo stability of the baseline contrasts}\label{sec:stability}
The baseline experiment was repeated under \RobSeeds{} independent master seeds. Periodic AHP ($k=10$) had lower regret than adaptive PSDF in \RobAHPTenAllNeg{} of \RobSeeds{} runs. The exponentially weighted ridge had lower regret than the rolling-window ridge in \RobEWLower{} of \RobSeeds{} runs (unadjusted $p<0.05$ in \RobEWSig; paired differences between \RobEWMin{} and \RobEWMax{}). Additive aggregation had lower regret than multiplicative aggregation in \RobAddLower{} runs but was Holm-significant in only \RobAddSig{} (differences between \RobAddMin{} and \RobAddMax{}). The AHP advantage and the exponential-weighting advantage are therefore stable in direction, whereas the multiplicative-versus-additive contrast is small and sits at the boundary of statistical detectability, so no conclusion is drawn about it beyond its small magnitude.

\section{Empirical Application: Monthly Provincial Tourism Data}\label{sec:empirical}
\subsection{Data}
The application uses the monthly provincial workbooks of the Ministry of Tourism and Sports for 2024 (revised), 2025 (revised) and 2026 (January--March revised, April--September preliminary) \citep{mots2026}. Each monthly sheet reports, for the 77 provinces (Bangkok included), the accommodation occupancy rate, accommodation guests, visitors (total, Thai, foreign) and tourism receipts (total, Thai, foreign), each with the corresponding value for the same month of the previous year. Merging the three workbooks yields a balanced panel of 77 provinces over 45 months (January 2023--September 2026). The overlapping vintages were compared: all 924 province-month records shared by the 2024 sheets of the 2024 and 2025 workbooks, and all 693 shared by the 2025 sheets of the 2025 and 2026 workbooks, were identical in every variable. Thai and foreign visitors and receipts sum to the totals in every record, no values are missing, the monthly provincial sums reproduce the published 2024 annual totals (353{,}556{,}836 recorded visits; THB 2{,}924{,}506.59 million), and the January 2025 provincial sums equal the published national total row. Two province names with differing spellings across workbooks were harmonized. April--September 2026 values are preliminary and may be revised. All of these checks are reproduced by the panel-construction script in the supplementary material.

\subsection{Design}
The design below was fixed before any results were computed, and no parameter was tuned. Every month $t$ from February 2024 to September 2026 (32 decision months) is a decision period with the 77 provinces as alternatives. The ranking at month $t$ uses only information measured in month $t-1$.

\paragraph{Features, risk and outcome.} Five benefit features measured in month $t-1$ are percentile-ranked across provinces within the month to lie in $[0,1]$: the occupancy rate, the logarithm of visitors, the foreign share of visitors, the logarithm of receipts per visitor, and accommodation guests per visitor. The risk variable is the percentile rank of the coefficient of variation of monthly visitors over the twelve months ending in $t-1$. The outcome for province $i$ in month $t$ is the within-month $z$-score of the year-on-year log growth of receipts, which removes seasonality and the common national shock. Training pairs for the decision at $t$ are $(X_{s-1},Y_s)$ for months $s\le t-1$, so each outcome enters only after it has been observed.

\paragraph{Methods.} All PSDF variants use the multiplicative score with $\lambda=1$, $\alpha=0.1$ and $N_{\min}=12$ records: rolling ridge with $L=12$ months (PSDF), expanding-window ridge, exponentially weighted ridge (0.9 per month), and weights frozen at the first eligible fit. Further comparators are uniform weights, CRITIC weights computed from the previous month's features, a signed rolling ridge forecast (the same ridge fit, ranking by the predicted outcome without non-negativity, normalization or the multiplicative risk transformation), a momentum rule (ranking by the previous month's outcome), and random selection (expected selected outcome 0 by construction of the $z$-score). No AHP comparator is included because no pairwise judgments are available.

\paragraph{Metrics and inference.} Each month the top ten provinces by score are selected. Regret is the mean outcome of the ex-post best ten provinces minus the mean outcome of the selected ten (top five as a secondary metric), and the mean Spearman correlation between score and outcome is reported. Differences are paired across the 32 months and intervals use a moving-block bootstrap (block length 3, 10{,}000 resamples), so they are approximate given the short, serially dependent series. The mean ex-post best-ten outcome is \EmpBestTen{} (top five: \EmpBestFive{}), which is the regret of random selection.

\subsection{Results}
\begin{table}[t]
\centering\small
\caption{Empirical application: selection regret in outcome $z$-score units (32 decision months, 77 provinces). Differences are relative to PSDF (rolling ridge); negative values favor the comparator. Intervals are moving-block bootstrap 95\% intervals.}\label{tab:empirical}
\begin{tabular}{lrrrl}
\toprule
Method & Top-10 regret & Top-5 regret & Spearman & Diff.\ vs PSDF (top 10) \\
\midrule
PSDF (rolling ridge, $L=12$) & 1.578 & 2.172 & $+0.007$ & -- \\
PSDF, expanding window & 1.600 & 2.227 & $+0.003$ & $+0.022$ [$-0.016$, $+0.080$] \\
PSDF, exponentially weighted & 1.605 & 2.160 & $+0.008$ & $+0.027$ [$-0.004$, $+0.072$] \\
PSDF, static weights & 1.661 & 2.338 & $-0.051$ & $+0.083$ [$-0.011$, $+0.209$] \\
Uniform weights & 1.674 & 2.356 & $-0.044$ & $+0.096$ [$+0.052$, $+0.154$] \\
CRITIC weights & 1.705 & 2.386 & $-0.046$ & $+0.126$ [$+0.070$, $+0.203$] \\
Signed rolling ridge forecast & 1.148 & 1.506 & $+0.164$ & $-0.430$ [$-0.638$, $-0.262$] \\
Momentum (previous month) & 0.620 & 0.812 & $+0.527$ & $-0.958$ [$-1.263$, $-0.616$] \\
Random selection (expected) & 1.513 & 2.026 & 0 & $-0.065$ [$-0.189$, $+0.072$] \\
\bottomrule
\end{tabular}

\end{table}
PSDF had lower regret than uniform and CRITIC weights, with intervals excluding zero, and its learned weights were concentrated on occupancy (mean \EmpWOcc) and the foreign share (\EmpWForeign), with \EmpWGuests, \EmpWReceipts{} and \EmpWVisitors{} on guests per visitor, receipts per visitor and visitors. However, its top-ten regret was not distinguishable from that of random selection (PSDF minus random: \EmpTenDiff, 95\% interval [\EmpTenLo, \EmpTenHi]), and at the top-five level it was slightly worse than random (\EmpFiveDiff, 95\% interval [\EmpFiveLo, \EmpFiveHi]). The mean rank correlation with next-month outcomes was close to zero (\EmpSpearman). The signed ridge forecast, which uses the same training data, and the momentum rule had substantially lower regret. PSDF's regret was \EmpFirstHalf{} in the first and \EmpSecondHalf{} in the second half of the evaluation period.

These results do not support a claim that PSDF, as specified, provides useful rankings in this application. A plausible reading is that the outcome is dominated by persistence and mean reversion that require signed coefficients or lagged outcomes, whereas PSDF represents only non-negative weights on benefit features; the results do not establish this mechanism. The application evaluates one outcome definition fixed in advance, ranks provinces by forecast growth and does not evaluate policy decisions or stated preferences; other objectives, such as the level of yield, were not examined.

\section{Discussion}
\subsection{Interpreting adaptive weighting}
PSDF uses outcome feedback to update criterion weights, unlike preference elicitation in which a decision maker directly expresses relative importance. A coefficient that predicts observed outcomes may reflect the outcome-generating process rather than normative preferences. The framework should therefore be read as an outcome-feedback-based decision-support mechanism, and its recovery of a known preference vector in simulation is conditional on the simulated outcome model.

\subsection{Separating weighting and aggregation}
The ablations indicate that the performance of PSDF relative to static comparators is attributable mainly to adaptive weighting and not to the multiplicative risk penalty: replacing adaptive weights with period-zero or initial-AHP weights increased regret by roughly 0.07--0.08, whereas replacing multiplicative with additive aggregation changed regret by only \AddDiff{} (Holm $p=\AddHolm$), a contrast whose significance is not stable across seeds. The fixed-weight comparator is estimated from a single period of feedback and is therefore a weak benchmark, which should be considered when interpreting the size of that gap.

\subsection{Interpreting the AHP comparison}
AHP and PSDF rely on different information channels. AHP receives pairwise judgments generated from the current preference vector at each elicitation round, whereas PSDF learns from historical outcome labels through rolling ridge regression. The comparison therefore reflects performance under different feedback assumptions and not equal information.

Two design features favor periodic AHP under the baseline. First, after the reversal AHP can update at the next scheduled round, whereas PSDF must wait for its rolling window to be dominated by post-reversal observations, which explains why PSDF regret is concentrated at $t=30$--44. Second, the baseline reversal at $t=30$ coincides with elicitation rounds for $k=5$ and $k=10$; with $k=20$, whose next round falls at $t=40$, AHP regret exceeded that of PSDF. The advantage of periodic AHP is thus sensitive to the alignment between elicitation schedule and change point.

The regime map shows that the ordering depends on a second factor: whether stated preferences coincide with the weights that govern outcomes. When they do not, AHP faithfully elicits the wrong target, whereas PSDF learns from outcomes. This argument is conditional on the gap magnitude, which is a modelling assumption. A fair assessment of information efficiency would also require measuring feedback acquisition cost, elicitation effort, and the value of the information each method receives.

\subsection{Implications for decision support}
The framework may be relevant to sequential decision problems in which outcomes are observed over time and the environment may change, particularly when repeated elicitation is costly or infeasible, or when stated and outcome-relevant preferences may diverge. Practical deployment requires evidence that outcome labels are valid, feedback is representative, and the learned model does not reinforce undesirable historical patterns. The simulation is a controlled methodological evaluation, not a substitute for field validation or user studies. The real-data application shows that the framework's non-negative benefit-weight structure can fail to capture the signal in observed outcomes, so applicability must be checked against simpler signed or persistence-based benchmarks.

\section{Limitations}
First, the simulation uses synthetic preferences, risk values, and outcomes, so results are conditional on the specified data-generating process; the semi-synthetic variant uses real provincial features for only three criteria. Second, AHP and PSDF use different information channels, and in the baseline the AHP comparator is given noisy but direct access to the current true preference vector; this asymmetry is structural and should not be read as evidence of the inherent superiority of either method. Third, the baseline reversal is aligned with the elicitation schedule for $k=5$ and $k=10$, and only a single abrupt reversal was examined; gradual drift and multiple reversals are not studied here. Fourth, the factorial regime map was specified after the baseline result was known, the preference-gap magnitudes are not calibrated to observed stated--revealed discrepancies, and the significance classification of cells with differences close to zero is not stable across random seeds. Fifth, interpreting regression coefficients as criterion weights depends on the relationship between observed outcomes and the decision objective. Sixth, the rolling-window estimator may be sensitive to abrupt changes, limited observations and misspecification, and an exponentially weighted alternative performed slightly better. Seventh, the multiplicative risk penalty imposes a particular functional form. Eighth, the fixed-weight comparator relies on a single period of feedback. Ninth, the simulation does not constitute empirical validation, and the real-data application covers one outcome definition, 32 decision months, no AHP comparison, and preliminary values for April--September 2026. Finally, software-level temporal-integrity tests cannot rule out all forms of bias, and independent code review and external replication remain advisable.

\section{Conclusion}
This paper evaluates a feedback-driven approach to adaptive multi-criteria decision support, PSDF, which combines rolling outcome-based estimation of benefit weights with explicit risk-adjusted aggregation, and asks when it has an advantage over AHP-based preference elicitation. Under the baseline configuration, in which simulated AHP judgments are generated from the current true preferences with small noise, periodically re-elicited AHP achieved lower expected regret than PSDF (mean paired difference, AHP minus PSDF, \DiffAHPTen, 95\% CI [\DiffAHPTenLo, \DiffAHPTenHi]), and this ordering held in all \RobSeeds{} repeated runs. The factorial analysis shows that the ordering depends on two conditions: whether the preference reversal coincides with a scheduled elicitation round, and whether stated preferences coincide with the weights that govern outcomes. When the reversal falls shortly after an elicitation round the baseline difference is close to zero, and when the preference gap is large PSDF achieved lower regret in both synthetic and semi-synthetic environments. Adaptive weighting accounted for most of PSDF's advantage over static, expanding-window, CRITIC and uniform weights; the form of risk aggregation changed regret only marginally; and an exponentially weighted ridge was marginally better than the rolling window.

In an application to monthly provincial tourism statistics (77 provinces, 32 decision months), PSDF improved on uniform and CRITIC weights but was not distinguishable from random selection at the top-ten level and was outperformed by a signed ridge forecast and a momentum rule. The simulation advantages reported above therefore did not carry over to this application.

These findings are conditional on a synthetic data-generating process, an exploratory factorial specified after the baseline result, and preference-gap magnitudes that have not been calibrated to observed stated--revealed discrepancies. Empirical validation on sequential outcome data with matched elicitation by decision makers, and a treatment of elicitation cost, are required before the results are extended to applied settings.

\section*{Data and Code Availability}
The simulation engine, experiment runner, temporal-integrity tests, seed-robustness script, the panel-construction script and the empirical-application script, with the commands that regenerate every table and number of Sections~\ref{sec:results} and~\ref{sec:empirical} (simulation run identifier \texttt{\RunID}), are reproduced in full in the supplementary material and archived at Zenodo (\url{https://doi.org/\RepoDOI}), together with the master seed, the machine-readable manifests and the merged province-by-month panel (\texttt{panel.csv}). The semi-synthetic environment uses provincial values from the publicly available Ministry of Tourism and Sports workbook cited in the references (the extracted values reconcile with the published national totals), and the empirical application uses the monthly provincial workbooks cited in the references, which are available from the Ministry of Tourism and Sports. \todo{Replace the identifier in \texttt{\textbackslash RepoDOI} (top of this file) by the DOI of the Zenodo release that contains the corrected files; the DOI of the earlier, incorrect upload must not be used.}

\section*{Declarations}
\noindent\textbf{Funding:} This research received no specific grant from any funding agency in the public, commercial, or not-for-profit sectors.\\
\textbf{Conflict of interest:} The author declares that she has no competing interests.\\
\textbf{Ethics approval:} Not applicable to the synthetic simulation study. The empirical application uses publicly available aggregate provincial statistics and does not involve human participants.\\
\textbf{CRediT author statement:} Patcharaporn Samarnggoon: Conceptualization, Methodology, Software, Formal analysis, Investigation, Data curation, Writing -- original draft, Writing -- review \& editing.\\
\textbf{Declaration of generative AI and AI-assisted technologies in the writing process:} During the preparation of this work, the author used AI assistance, including Claude (Anthropic), for drafting and coding support. After using this tool, the author reviewed and edited the content as needed and takes full responsibility for the content of the published article.

\begin{thebibliography}{99}\small
\bibitem[Ayan et~al.(2023)]{ayan2023} Ayan, B., Abac{\i}o\u{g}lu, S., \& Basilio, M.~P. (2023). A comprehensive review of the novel weighting methods for multi-criteria decision-making. \emph{Information}, 14(5), 285. \url{https://doi.org/10.3390/info14050285}
\bibitem[Bengs et~al.(2021)]{bengs2021} Bengs, V., Busa-Fekete, R., El~Mesaoudi-Paul, A., \& H\"ullermeier, E. (2021). Preference-based online learning with dueling bandits: A survey. \emph{Journal of Machine Learning Research}, 22.
\bibitem[Cinelli et~al.(2020)]{cinelli2020} Cinelli, M., Kadzi\'nski, M., Gonzalez, M., \& S{\l}owi\'nski, R. (2020). How to support the application of multiple criteria decision analysis? Let us start with a comprehensive taxonomy. \emph{Omega}, 96, 102261. \url{https://doi.org/10.1016/j.omega.2020.102261}
\bibitem[H\"ullermeier and S{\l}owi\'nski(2024a)]{hullermeier2024a} H\"ullermeier, E., \& S{\l}owi\'nski, R. (2024a). Preference learning and multiple criteria decision aiding: Differences, commonalities, and synergies---part I. \emph{4OR}, 22, 179--209. \url{https://doi.org/10.1007/s10288-023-00560-6}
\bibitem[H\"ullermeier and S{\l}owi\'nski(2024b)]{hullermeier2024b} H\"ullermeier, E., \& S{\l}owi\'nski, R. (2024b). Preference learning and multiple criteria decision aiding: Differences, commonalities, and synergies---part II. \emph{4OR}, 22, 313--349. \url{https://doi.org/10.1007/s10288-023-00561-5}
\bibitem[Lu and Zhang(2026)]{lu2026} Lu, C., \& Zhang, Y. (2026). A dynamic hybrid weighting framework for teaching effectiveness evaluation in multi-criteria decision-making: Integrating interval-valued intuitionistic fuzzy AHP and entropy triggering. \emph{Entropy}, 28(2), 241. \url{https://doi.org/10.3390/e28020241}
\bibitem[Ministry of Tourism and Sports(2025)]{mots2025} Ministry of Tourism and Sports (2025). \emph{Provincial Tourism Statistics, January--December 2024: Revised 28 August 2025}. Bangkok: Ministry of Tourism and Sports. \url{https://secretary.mots.go.th/tourism-stat}
\bibitem[Ministry of Tourism and Sports(2026)]{mots2026} Ministry of Tourism and Sports (2026). \emph{Domestic tourism statistics by province, monthly, 2024--2026 (revised and preliminary workbooks)}. Bangkok: Ministry of Tourism and Sports. \url{https://www.mots.go.th/news/category/411}
\bibitem[Saaty(1980)]{saaty1980} Saaty, T.~L. (1980). \emph{The Analytic Hierarchy Process}. McGraw-Hill.
\bibitem[Yang et~al.(2025)]{yang2025} Yang, S., Liao, H., \& Wu, X. (2025). Prescriptive analytics for dynamic multi-criterion decision making considering learned knowledge of alternatives. \emph{Expert Systems with Applications}. \url{https://doi.org/10.1016/j.eswa.2024.126350}
\end{thebibliography}
\end{document}
